\documentclass[11pt,a4paper]{article}

% ---------- packages ----------
\usepackage[utf8]{inputenc}
\usepackage[T1]{fontenc}
\usepackage[english]{babel}
\usepackage{microtype}
\usepackage[margin=2.5cm]{geometry}
\usepackage{amsmath,amssymb,amsthm}
\usepackage{booktabs}
\usepackage{array}
\usepackage{enumitem}
\usepackage{xcolor}
\usepackage[hidelinks]{hyperref}
\usepackage{titlesec}
\usepackage{parskip}
\usepackage{eurosym}

% ---------- fine-tuning ----------
\setlength{\parindent}{0pt}
\titleformat*{\section}{\Large\bfseries}
\titleformat*{\subsection}{\large\bfseries}
\titleformat*{\subsubsection}{\normalsize\bfseries}

\newcommand{\code}[1]{\allowbreak\texttt{\small #1}\allowbreak}

% allow slightly loose typesetting for long inline monospace identifiers
\emergencystretch=3em

% ---------- meta ----------
\title{\textbf{A3I Project --- Final Report}\\[0.3em]
\large Day-Ahead Price Prediction and Battery Arbitrage\\
on the IT\_NORD Bidding Zone}
\author{}
\date{}

\begin{document}
\maketitle

% ---------- 1. Executive Summary ----------
\section{Executive Summary}

Italian electricity prices on the IT\_NORD day-ahead market move with a
combination of meteorological drivers (solar radiation, wind, temperature),
calendar effects (daily and weekly demand cycles) and autoregressive
inertia. When a battery-storage operator has even a reasonable guess of
tomorrow's 24 hourly prices, they can turn that forecast into cash by
charging at the low hours and discharging at the high ones --- as long as
the spread covers the physical wear-and-tear of a charge/discharge cycle,
known in the industry as the Levelized Cost of Storage (LCOS).

This project delivers an end-to-end pipeline that goes from the raw
public data (ENTSO-E prices, Open-Meteo weather) all the way to an
operational trading policy for a 1\,MWh battery. The final policy is
trained with \emph{Decision-Focused Learning} (DFL), which directly
optimizes the quantity that matters --- the net profit after LCOS ---
rather than the statistical accuracy of the price forecast.

Three technical contributions hold the project together:

\begin{itemize}[leftmargin=1.5em]
  \item \textbf{An autoencoder trained on the quiet 2019--2020 market}
    identifies the 2021--2022 regime shift as a sustained spike in
    reconstruction error. No labels and no hand-picked crisis dates are
    needed; the anomaly falls out of the data.

  \item \textbf{An XGBoost price model}, evaluated on a genuine
    2024--2025 hold-out, is analyzed with permutation importance, SHAP
    (global and local) and Boruta's hypothesis-testing procedure. The
    analysis confirms that the model has internalized the
    \emph{merit-order effect} --- renewable generation cannibalising
    wholesale prices --- and lets us defend the claim that each retained
    feature carries genuine information rather than a spurious
    correlation.

  \item \textbf{A DFL battery-arbitrage policy} trained with a
    softmax-relaxed regret loss that folds LCOS into the training
    objective. Compared to the standard predict-then-optimize baseline
    on the same 2024--2025 hold-out, the DFL policy delivers a higher
    net profit while making fewer trades, i.e.\ it absorbs the LCOS more
    effectively.
\end{itemize}


% ---------- 2. Problem Setting ----------
\section{Problem Setting}

\subsection{The market}
The Italian day-ahead market clears one price per hour, per bidding
zone. IT\_NORD covers the industrial heartland of the country and is the
largest zone by traded volume. Wholesale clearing follows the standard
\emph{merit-order}: the cheapest generators (typically renewable, with
near-zero marginal cost) are dispatched first, and the price settles at
the marginal cost of the last plant called into service. The mechanism
implies two properties that our model will need to learn:

\begin{itemize}[leftmargin=1.5em]
  \item \textbf{Strong autoregressive persistence.} Industrial load
    cycles are sticky --- today's 18:00 price looks a lot like
    yesterday's 18:00 price. A shock propagates for days.
  \item \textbf{Asymmetric price response to weather.} A sunny day in
    July depresses mid-day prices because solar cannibalizes the
    dispatch stack; the same radiation level at 03:00 has no effect
    because solar is dark. The two inputs \emph{interact}.
\end{itemize}

\subsection{The operator's problem}
For a battery operator, the day-ahead market defines a daily arbitrage
game. Given a point forecast $\hat{y}\in\mathbb{R}^{24}$ of tomorrow's
24 hourly prices, the operator chooses one charge hour $h_c$ and one
discharge hour $h_d$. The realized profit is
\[
\pi = y_{h_d} - y_{h_c} - C_{deg}
\]
if we decide to trade, and $0$ if we stay idle (the LCOS penalty is only
paid when the battery actually cycles). The \emph{optimal} decision is
to trade whenever the realized spread covers the LCOS:
\[
\pi^{*} \;=\; \max\!\left(0,\;\; \max_{h} y_{h} - \min_{h} y_{h} - C_{deg}\right).
\]
The gap between $\pi$ (what we actually earn, using the forecast) and
$\pi^{*}$ (what an oracle with perfect foresight would earn) is the
\emph{regret}. Classical ML minimizes the forecast error; our pipeline
minimizes the regret directly.

\subsection{Why this is a two-stage problem}
The course lecture \code{02\_\_Limitations\_and\_Opportunities.ipynb}
makes a sharp point: \textbf{for purely linear cost functions with
uncertain coefficients, a well-specified predict-then-optimize model
already converges to the optimum --- so DFL offers little.} Where DFL
genuinely pays off is in \emph{two-stage stochastic optimization}, where
the cost is non-linear in the uncertain parameter because of recourse
actions (here: the idle option, which is a discrete non-linearity in
$y$).

Battery arbitrage with LCOS is precisely that. Our DFL formulation folds
$C_{deg}$ into the training loss, so the network learns to produce
forecasts whose implied trading decision is decision-relevant, not just
statistically close.


% ---------- 3. Data and Features ----------
\section{Data and Features}

\subsection{Sources}
\begin{center}
\begin{tabular}{@{}lllc@{}}
\toprule
\textbf{Source} & \textbf{Variables} & \textbf{Period} & \textbf{Resolution} \\
\midrule
ENTSO-E    & Day-ahead price IT\_NORD (EUR/MWh)               & 2019--2025 & Hourly \\
Open-Meteo & Temperature (2 m), direct radiation, wind (10 m) & 2019--2025 & Hourly \\
\bottomrule
\end{tabular}
\end{center}
Both sources are aligned strictly on UTC; Rome-local timestamps would
have introduced DST artifacts that propagate into the daily-cycle
features.

\subsection{Engineered features}
The feature set passed to XGBoost is summarized in
Table~\ref{tab:features}.

\begin{table}[h]
\centering
\small
\begin{tabular}{@{}>{\raggedright\arraybackslash}p{3cm}>{\raggedright\arraybackslash}p{6cm}>{\raggedright\arraybackslash}p{6cm}@{}}
\toprule
\textbf{Group} & \textbf{Features} & \textbf{Rationale} \\
\midrule
Weather & Temperature\_C, \code{Solar\_Radiation\_W}, \code{Wind\_Speed\_ms}
              & Merit-order + thermal demand \\
Calendar (raw)& {Is\_Weekend} & Industrial demand drop on Sat--Sun \\
Calendar (cyc)& {Hour\_Sin}, \code{Hour\_Cos}, \code{Month\_Sin}, \code{Month\_Cos}
              & Continuous wrap-around at midnight and end-of-year \\
Autoregressive & {Price\_Lag\_24h}, \code{Price\_Lag\_168h}, \code{Price\_Rolling\_Mean\_24h}
              & Daily + weekly persistence; rolling on the lag, not target \\
\bottomrule
\end{tabular}
\caption{Feature groups passed to the XGBoost price model.}
\label{tab:features}
\end{table}

The rolling mean is computed on \code{Price\_Lag\_24h} rather than on
the target itself --- the standard trick to avoid look-ahead bias.
Cyclical encoding of hour and month matters for tree-based learners
because it puts hour~23 next to hour~0 on the unit circle, which ordinal
encoding does not.


% ---------- 4. Autoencoder ----------
\section{Anomaly Detection via Autoencoder}

A shallow autoencoder ($5\!\to\!16\!\to\!8\!\to\!3\!\to\!8\!\to\!16\!\to\!5$)
is trained on the quiet years 2019--2020 with the five-dimensional
state vector (price, three weather features, \code{Price\_Lag\_24h}) as
both input and target. The bottleneck forces the network to learn only
the ``normal'' manifold.

At inference time we run the trained encoder-decoder on the full
2019--2025 dataset and compute the per-row reconstruction MSE. A
seven-day rolling mean of the MSE gives a smoothed \emph{anomaly score}.
A threshold set at the 99th percentile of the training-period MSE marks
any sustained excursion above it as a structural anomaly.

The score stays flat through 2019--2020, climbs through late 2021, and
peaks in autumn 2022 --- the point in time when European gas prices,
driven by the supply shock, reset the entire wholesale electricity
stack. The model picks this up \emph{without} ever having been told
about the gas crisis; the anomaly falls out of the reconstruction error
alone.

\paragraph{Operational use.} An operator could run this score live: when
it crosses the threshold, retrain or flag the downstream forecasting
model for close monitoring, since it was trained under an implicit
stationarity assumption that the market is about to break.


% ---------- 5. XGBoost ----------
\section{Price Prediction with XGBoost}

\subsection{Set-up}
A chronological split is used: 2019--2023 for training (with the last
10\% carved off for early stopping), 2024--2025 for the hold-out.
Random splits would contaminate the training set with future
information and inflate the reported metrics. The model is an
\code{XGBRegressor} with 500 trees (early-stopped), learning rate
$0.05$, max depth $6$.

On the 2024--2025 hold-out, the model obtains an MAE of roughly 12--15
\euro/MWh depending on the run (the exact number is reproducible with the
fixed seed in the notebook). That is materially better than a na\"ive
persistence baseline (predict today = yesterday), which hovers around
25 \euro/MWh MAE on the same window.

\subsection{Feature selection analysis}
Three complementary techniques, drawn from the course notebooks
\code{3\_\_Using\_a\_Non-Linear\_Model}, \code{4\_\_Additive\_Feature\_Attribution}
and \code{5\_\_All\_Relevant\_Feature\_Selection}, are applied in the
\code{Feature\_Selection\_Module.ipynb} companion notebook.

\paragraph{Permutation importance on the test set.}
Computed with \code{permutation\_importance} from
\code{sklearn.inspection}, $n=30$ repeats, scored in RMSE. The ranking
that comes out is dominated by the autoregressive features
(\code{Price\_Lag\_24h}, \code{Price\_Rolling\_Mean\_24h},
\code{Price\_Lag\_168h}) followed by \code{Solar\_Radiation\_W} and the
cyclical hour features, with \code{Is\_Weekend} trailing. Error bars
give us a principled way to identify features whose importance cannot
be distinguished from zero.

\paragraph{SHAP, global and local.}
The mean-$|$SHAP$|$ bar plot agrees with the permutation ranking. The
beeswarm plot adds sign information: high lag prices push the forecast
up monotonically (as expected for persistence), while high solar
radiation pushes it down. The dependency plot of
\code{Solar\_Radiation\_W} coloured by \code{Hour\_Sin} is the most
informative single figure in the whole analysis: the negative
contribution of solar is strongly modulated by the hour of day, with
the effect concentrated in daylight hours. This is the merit-order
effect --- recovered automatically by the model from raw inputs,
without the modeller hand-coding it.

\paragraph{Boruta.}
BorutaPy is run with the XGBoost backbone, 50 iterations. Each real
feature is tested against its \emph{shadow} (a shuffled copy of itself)
under the null hypothesis that its importance is no greater than what
chance alone would produce. The procedure returns three buckets:
\textbf{confirmed relevant}, \textbf{tentative}, and \textbf{rejected}.
On our data, the autoregressive features and the
weather-plus-cyclical-time bundle are confirmed; \code{Is\_Weekend}
typically lands in the tentative band, which is consistent with its
signal being already partially absorbed by \code{Price\_Lag\_168h} ---
the mediator / redundancy phenomenon described in the last lecture of the Explainable block.

\paragraph{Minimal-set sanity check.}
Starting from the permutation-importance ranking, we retrain the
XGBoost on progressively smaller subsets. The test RMSE drops fast in
the first 3--5 features (lags + rolling mean + solar + one cyclical
hour), then plateaus. A deployment-oriented model could drop the
weakest 2--3 features without measurable accuracy loss, which would
halve feature-engineering cost at inference time.

\subsection{What feature selection bought us}
Two concrete things:

\begin{itemize}[leftmargin=1.5em]
  \item \textbf{Defensibility.} The XGBoost model is not just accurate,
    it is \emph{interpretable}: we can point at each retained feature
    and present either a SHAP plot or a Boruta hypothesis test that
    justifies keeping it. This matters the moment a trading desk or a
    regulator asks ``why did your model recommend that trade?''.

  \item \textbf{Cleaner downstream inputs.} The DFL battery-net in
    Section~\ref{sec:dfl} takes the 24 lagged prices as inputs; the
    feature-selection analysis tells us the rolling mean adds little
    beyond the raw lags, which guides any future extension of the DFL
    input set.
\end{itemize}


% ---------- 6. DFL ----------
\section{Decision-Focused Learning for Battery Arbitrage}\label{sec:dfl}

\subsection{Architecture}
A single-hidden-layer MLP ($24\!\to\!64\!\to\!24$) takes yesterday's 24
hourly prices (standardized with training-only statistics) and outputs
tomorrow's 24 price forecast. Two copies of the same architecture are
trained with different losses.

\begin{itemize}[leftmargin=1.5em]
  \item \textbf{PFL} --- plain MSE between predicted and true hourly
    prices. The LCOS idle rule is applied only at evaluation time, by
    checking whether the \emph{predicted} spread exceeds $C_{deg}$.
    This is what PFL \emph{is} by construction.

  \item \textbf{DFL} --- a softmax-relaxed regret with LCOS folded into
    the training loss. Writing
    $p^{dis}_{h}=\operatorname{softmax}(\hat y)_{h}$ and
    $p^{ch}_{h}=\operatorname{softmax}(-\hat y)_{h}$:
    \[
    \mathcal{L}_{\mathrm{DFL}}
      = \mathbb{E}\!\left[
          \pi^{*}\;-\;\operatorname{softplus}\!\Bigl(
            \textstyle\sum_{h}(p^{dis}_{h}-p^{ch}_{h})\,y_{h}\;-\;C_{deg}
          \Bigr)
        \right]
      + \lambda\,\mathbb{E}\!\left[\textstyle\sum_{h}p^{dis}_{h}\,p^{ch}_{h}\right].
    \]
    The first term is the oracle profit minus the soft-ReLU of the
    expected spread net of LCOS --- making the idle option a
    \emph{differentiable} choice. The second is a small collision
    penalty that discourages the model from placing discharge mass and
    charge mass on the same hour.
\end{itemize}

\subsection{Relation to the course's canonical DFL losses}
The lecture notes discuss three families of differentiable surrogates:
SPO+, self-contrastive, and softmax / perturbed-optimizer relaxations.
Our loss sits in the third family. It is \textbf{not} the Score Function
Gradient Estimator: SFGE (REINFORCE) would require sampling a discrete
action and using the log-derivative trick, whereas here the expectation
is computed in closed form via the softmax weights.

\subsection{Evaluation}
The daily tensor is split chronologically: days from 2019 through 2023
go to training, days from 2024 through 2025 go to the hold-out. Both
models train for 250 epochs with Adam at lr~$=10^{-2}$. At evaluation
time the discrete $\operatorname{argmax}$~/~$\operatorname{argmin}$ are
recovered from the point forecasts, the LCOS idle rule is applied to
both models, and profits are summed over the test days in euros (1~MWh
nominal battery).

On the hold-out we observe:

\begin{itemize}[leftmargin=1.5em]
  \item \textbf{PFL} captures the clear-cut days where the spread far
    exceeds LCOS, but frequently takes trades in marginal conditions
    where the realized spread falls just short --- its forecast
    regresses to the mean and underestimates volatility.
  \item \textbf{DFL} converges to a policy that is simultaneously
    \emph{more selective} (it trades less often on flat days) and
    \emph{more aggressive} (it captures the largest spreads more
    reliably). The net result is a higher cumulative profit over the
    test horizon, despite trading fewer days in total.
\end{itemize}

The cumulative-profit plot in the notebook shows the DFL curve
consistently above the PFL curve, with both well below the oracle ---
as expected, since the oracle has perfect foresight.

\subsection{Sensitivity and optimization-protocol analysis}
To probe the sensitivity of the DFL-vs-PFL comparison to both $C_{deg}$ and the neural-network random initialization, we ran four optimization protocols across $C_{deg} \in \{5, 10, 15, 20, 25, 30, 40\}$ \euro/MWh with five seeds per configuration: cold-start DFL (random init, 250 epochs at lr $10^{-2}$), warm-start DFL (initialized from the trained PFL weights, refined at lr $10^{-3}$), and two hybrid protocols that anneal the training loss linearly from pure MSE to pure regret over 400 epochs, with and without an amplification factor on the regret term.

Cold-start DFL dominates PFL on average, with the advantage peaking at roughly \euro5k in the central $C_{deg}$ range, but with seed-to-seed standard deviation of the same order as the mean. The three MSE-anchored protocols (warm-start, hybrid, amplified hybrid) all collapse to PFL for $C_{deg} \leq 25$
\euro/MWh — not because the networks fail to differ from PFL (we measured per-hour prediction differences of up to \euro 8/MWh) but because they make identical charge/discharge hour choices on nearly every test day. Increasing the regret term's weight by a factor of 20 in the hybrid protocol did not change this, indicating that the PFL optimum is a decision-local minimum from which small regret gradients cannot escape, regardless of their magnitude. At $C_{deg} \geq 30$
\euro/MWh all three MSE-anchored protocols recover a small ($\sim$\euro500) but stable advantage, the range in which the idle-vs-trade decision begins to dominate and the regret landscape diverges measurably from MSE's.

The combined picture is that on this problem the regret loss landscape is shaped by the discrete ordering of hourly prices, which is easy to get right on most days; MSE-basin solutions are therefore decision-optimal for most days even when their values are imperfect, and genuine DFL-driven decision changes require escape from the MSE basin — which costs seed variance. A deployed DFL pipeline on day-ahead arbitrage would accordingly either ensemble multiple cold-start runs to average over that variance, or restrict DFL to the high-$C_{deg}$ regime where warm-started refinement yields a modest but stable improvement. This behavior echoes the theoretical observation in the first lecture about DFL that regret is piecewise-constant by construction and admits multiple local minima.

% ---------- 7. Limitations ----------
\section{Limitations and Honest Caveats}

\begin{itemize}[leftmargin=1.5em]
  \item \textbf{Battery model is minimal.} We model a 1\,MWh /
    1~cycle-per-day battery with a scalar LCOS and no state of charge,
    efficiency losses, or ramp constraints. A production model would be
    a more involved mixed-integer linear program; the DFL
    \emph{principle} carries over, but the training loop would need a
    differentiable solver or a black-box surrogate for the MILP.

  \item \textbf{Single bidding zone.} The model is trained and
    evaluated on IT\_NORD only. Cross-zonal arbitrage and transmission
    constraints are out of scope.

  \item \textbf{No intra-day or balancing markets.} The real trading
    stack for a battery operator includes the intra-day continuous
    market and the ancillary-services / balancing markets, often with
    higher volatilities than day-ahead. Our pipeline only addresses the
    day-ahead leg, which is typically the base layer of an operator's
    strategy but not the whole of it.

  \item \textbf{Regime-shift robustness.} We flagged the 2021--2022
    crisis with the autoencoder, but we did not retrain separate models
    per regime. A deployed system should at minimum re-fit on rolling
    windows and treat autoencoder alarms as triggers for retraining.
\end{itemize}


% ---------- 8. Usefulness ----------
\section{What This Project Is Useful For}

The pipeline produces four artefacts that are useful in distinct
operational contexts.

\paragraph{1. A battery-arbitrage decision-support tool.}
For an operator running a grid-scale battery on the IT\_NORD market,
the DFL policy gives a daily charge-hour / discharge-hour recommendation
with an explicit idle option when the expected spread does not cover
LCOS. The model runs on free public data (ENTSO-E + Open-Meteo) and
fits in under ten seconds of inference per day on a laptop; it is
immediately deployable as a morning briefing, as a first draft for a
human trader to approve, or as the inner loop of an automated bidder.

\paragraph{2. A market-monitoring early-warning system.}
The autoencoder anomaly score, run live on the rolling 7-day window,
gives a single number that summarizes how far the market currently is
from its ``normal'' 2019--2020 manifold. A trading desk can use it as a
sanity check on any short-horizon forecast: a forecast produced during
a reconstruction-error spike deserves extra scepticism, and any
automated policy trained on pre-shock data should be re-examined or
paused.

\paragraph{3. A transparent, regulator-ready explanation layer.}
The XGBoost price model, combined with the SHAP dependency plots and
the Boruta relevance report, constitutes an explainability bundle that
can be submitted to a regulator or an internal compliance team. The
bundle answers both global (``which features matter?'') and local
(``why did the model recommend this trade, on this day?'') questions,
using the same mathematical primitives (Shapley values, hypothesis
testing) that are acceptable in both academic and regulated settings.

\paragraph{4. A pedagogical template for the
predict-then-optimize vs decision-focused question.}
The side-by-side PFL~/~DFL training loop, run on a real market problem
rather than a synthetic toy, demonstrates concretely the point made in
\code{02\_\_Limitations\_and\_Opportunities.ipynb}: DFL earns its keep
in two-stage problems with recourse actions, and here the idle option
is exactly that kind of non-linearity.

\bigskip
Beyond the narrow use case of battery arbitrage, the same pipeline
skeleton --- public data ingestion, cyclical + autoregressive feature
engineering, autoencoder-based regime detection, explainable tree
model, DFL for a downstream decision --- ports cleanly to adjacent
problems: demand-response scheduling, EV-fleet charging optimization,
district-heating dispatch, or industrial peak shaving. All of these
share the same structure (predict an hourly profile $\rightarrow$ take
a discrete decision at the low/high of the profile, with a fixed cost
for acting), and therefore benefit from the same DFL treatment.


% ---------- 9. Deliverables ----------
\section{Deliverables}

\begin{itemize}[leftmargin=1.5em]
  \item \code{main\_restructured.ipynb} --- the end-to-end pipeline
    from data ingestion to the DFL evaluation.
  \item \code{Final\_Report.tex} (this document) --- the write-up of the project, with a structured presentation of the
    problem, the methods, the results, and the limitations.
\end{itemize}


% ---------- 10. Closing ----------
\section{Closing Remarks}

The project demonstrates that aligning the machine-learning loss with
the downstream operational utility --- regret minimization under LCOS
--- yields an honestly better trading policy than a standard MSE-driven
forecaster, on held-out days and under a reasonable cost assumption.
The accompanying feature-selection analysis makes the price model
itself interpretable and defensible, with mean-$|$SHAP$|$, permutation
importance, and Boruta hypothesis tests all pointing at the same core
drivers: autoregressive persistence, solar radiation with its clear
interaction with the hour of day, and the weekly demand cycle.

The pipeline is intentionally minimal --- one bidding zone, one
battery, one cycle per day --- because the goal was methodological
clarity. Scaling it up to a real operator's stack is an engineering
exercise; the ML and DFL foundations shown here are the non-obvious
part, and they generalize.

\end{document}